% GENERATED by scripts/generate-mindmaps.mjs — do not hand-edit.
\begin{tikzpicture}[
  mindmap,
  every node/.style={
    concept,
    execute at begin node=\hskip0pt,
    align=center,
    font=\sffamily
  },
  mm root/.style={
    concept color=graphite,
    text=ghostwhite,
    font=\sffamily\bfseries,
    minimum size=3.35cm,
    text width=2.75cm
  },
  mm section/.style={
    font=\sffamily\bfseries\scriptsize,
    minimum size=2.15cm,
    text width=1.68cm
  },
  mm topic/.style={
    font=\sffamily\tiny,
    minimum size=1.55cm,
    text width=1.08cm
  }
]
  \node[mm root] (root) at (0,0) {Chapter 7\\Reactive Systems, SMV, and Explicit Model Checking};
  \node[mm section, concept color=brightsky] (s1) at (80.27:6.93cm) {Why Reactive Systems?};
  \node[mm topic, concept color=brightsky!44!white] (t1_1) at (85.14:13.32cm) {Closed vs. Open Systems};
  \node[mm topic, concept color=brightsky!44!white] (t1_2) at (75.41:13.32cm) {The cost of failure};
  \node[mm section, concept color=vvteal] (s2) at (41.35:6.93cm) {The Model Checking Problem};
  \node[mm topic, concept color=vvteal!44!white] (t2_1) at (65.68:13.32cm) {What model checking is};
  \node[mm topic, concept color=vvteal!44!white] (t2_2) at (55.95:13.32cm) {The quality of model and specification};
  \node[mm topic, concept color=vvteal!44!white] (t2_3) at (46.22:13.32cm) {Universal and existential model checking};
  \node[mm topic, concept color=vvteal!44!white] (t2_4) at (36.49:13.32cm) {A brief historical note};
  \node[mm topic, concept color=vvteal!44!white] (t2_5) at (26.76:13.32cm) {Verification vs. validation};
  \node[mm topic, concept color=vvteal!44!white] (t2_6) at (17.03:13.32cm) {The state-explosion problem};
  \node[mm section, concept color=vvviolet] (s3) at (2.43:6.93cm) {Kripke Structures};
  \node[mm topic, concept color=vvviolet!44!white] (t3_1) at (7.3:13.32cm) {Kripke structure / transition system};
  \node[mm topic, concept color=vvviolet!44!white] (t3_2) at (-2.43:13.32cm) {Arbiter Kripke structure};
  \node[mm section, concept color=vvorange] (s4) at (-12.16:6.93cm) {Specifications: Safety and Liveness};
  \node[mm topic, concept color=vvorange!44!white] (t4_1) at (-12.16:13.32cm) {Two kinds of ``always good''};
  \node[mm section, concept color=tomato] (s5) at (-26.76:6.93cm) {Counterexamples: From Bugs to Witnesses};
  \node[mm topic, concept color=tomato!44!white] (t5_1) at (-21.89:13.32cm) {Finite prefixes for safety violations};
  \node[mm topic, concept color=tomato!44!white] (t5_2) at (-31.62:13.32cm) {Lassos for liveness violations};
  \node[mm section, concept color=brightsky] (s6) at (-85.14:6.93cm) {The SMV Modelling Language};
  \node[mm topic, concept color=brightsky!44!white] (t6_1) at (-41.35:13.32cm) {State variables and types};
  \node[mm topic, concept color=brightsky!44!white] (t6_2) at (-51.08:13.32cm) {The ASSIGN block};
  \node[mm topic, concept color=brightsky!44!white] (t6_3) at (-60.81:13.32cm) {Expression syntax};
  \node[mm topic, concept color=brightsky!44!white] (t6_4) at (-70.54:13.32cm) {The case expression};
  \node[mm topic, concept color=brightsky!44!white] (t6_5) at (-80.27:13.32cm) {Formal semantics of SMV};
  \node[mm topic, concept color=brightsky!44!white] (t6_6) at (-90:13.32cm) {Adding a self-loop: implicit modelling with TRANS};
  \node[mm topic, concept color=brightsky!44!white] (t6_7) at (-99.73:13.32cm) {Integer-range example: modulo-4 counter};
  \node[mm topic, concept color=brightsky!44!white] (t6_8) at (-109.46:13.32cm) {Non-determinism in SMV};
  \node[mm topic, concept color=brightsky!44!white] (t6_9) at (-119.19:13.32cm) {Input variables};
  \node[mm topic, concept color=brightsky!44!white] (t6_10) at (-128.92:13.32cm) {The define macro};
  \node[mm section, concept color=vvgray] (s7) at (-143.51:6.93cm) {Simulation with NuXMV};
  \node[mm topic, concept color=vvgray!44!white] (t7_1) at (-138.65:13.32cm) {Simulating the flip-flop};
  \node[mm topic, concept color=vvgray!44!white] (t7_2) at (-148.38:13.32cm) {Simulating the modulo-4 counter};
  \node[mm section, concept color=brightsky] (s8) at (-162.97:6.93cm) {Modules: Hierarchical Modelling};
  \node[mm topic, concept color=brightsky!44!white] (t8_1) at (-158.11:13.32cm) {The mechanical decimal counter};
  \node[mm topic, concept color=brightsky!44!white] (t8_2) at (-167.84:13.32cm) {Properties of the mechanical counter};
  \node[mm section, concept color=vvteal] (s9) at (-182.43:6.93cm) {Asynchronous Composition};
  \node[mm topic, concept color=vvteal!44!white] (t9_1) at (-177.57:13.32cm) {Synchronous vs. asynchronous};
  \node[mm topic, concept color=vvteal!44!white] (t9_2) at (-187.3:13.32cm) {Producer--consumer example};
  \node[mm section, concept color=vvviolet] (s10) at (-211.62:6.93cm) {Fairness Constraints};
  \node[mm topic, concept color=vvviolet!44!white] (t10_1) at (-197.03:13.32cm) {Temporal hierarchy: four key concepts};
  \node[mm topic, concept color=vvviolet!44!white] (t10_2) at (-206.76:13.32cm) {Formal definitions};
  \node[mm topic, concept color=vvviolet!44!white] (t10_3) at (-216.49:13.32cm) {Encoding fairness in NuSMV};
  \node[mm topic, concept color=vvviolet!44!white] (t10_4) at (-226.22:13.32cm) {Fair paths and model checking under fairness};
  \node[mm section, concept color=vvorange] (s11) at (-235.95:6.93cm) {LTL, CTL, and CTL*: an Overview};
  \node[mm section, concept color=tomato] (s12) at (-255.41:6.93cm) {Formal Methods in Context};
  \node[mm topic, concept color=tomato!44!white] (t12_1) at (-245.68:13.32cm) {Three categories of formal methods};
  \node[mm topic, concept color=tomato!44!white] (t12_2) at (-255.41:13.32cm) {Static analysis and abstract interpretation};
  \node[mm topic, concept color=tomato!44!white] (t12_3) at (-265.14:13.32cm) {The verification--synthesis spectrum};

  \begin{scope}[on background layer, line cap=round]
    \draw[brightsky, line width=11pt] (root) -- (s1);
    \draw[brightsky!60!white, line width=5.5pt] (s1) -- (t1_1);
    \draw[brightsky!60!white, line width=5.5pt] (s1) -- (t1_2);
    \draw[vvteal, line width=11pt] (root) -- (s2);
    \draw[vvteal!60!white, line width=5.5pt] (s2) -- (t2_1);
    \draw[vvteal!60!white, line width=5.5pt] (s2) -- (t2_2);
    \draw[vvteal!60!white, line width=5.5pt] (s2) -- (t2_3);
    \draw[vvteal!60!white, line width=5.5pt] (s2) -- (t2_4);
    \draw[vvteal!60!white, line width=5.5pt] (s2) -- (t2_5);
    \draw[vvteal!60!white, line width=5.5pt] (s2) -- (t2_6);
    \draw[vvviolet, line width=11pt] (root) -- (s3);
    \draw[vvviolet!60!white, line width=5.5pt] (s3) -- (t3_1);
    \draw[vvviolet!60!white, line width=5.5pt] (s3) -- (t3_2);
    \draw[vvorange, line width=11pt] (root) -- (s4);
    \draw[vvorange!60!white, line width=5.5pt] (s4) -- (t4_1);
    \draw[tomato, line width=11pt] (root) -- (s5);
    \draw[tomato!60!white, line width=5.5pt] (s5) -- (t5_1);
    \draw[tomato!60!white, line width=5.5pt] (s5) -- (t5_2);
    \draw[brightsky, line width=11pt] (root) -- (s6);
    \draw[brightsky!60!white, line width=5.5pt] (s6) -- (t6_1);
    \draw[brightsky!60!white, line width=5.5pt] (s6) -- (t6_2);
    \draw[brightsky!60!white, line width=5.5pt] (s6) -- (t6_3);
    \draw[brightsky!60!white, line width=5.5pt] (s6) -- (t6_4);
    \draw[brightsky!60!white, line width=5.5pt] (s6) -- (t6_5);
    \draw[brightsky!60!white, line width=5.5pt] (s6) -- (t6_6);
    \draw[brightsky!60!white, line width=5.5pt] (s6) -- (t6_7);
    \draw[brightsky!60!white, line width=5.5pt] (s6) -- (t6_8);
    \draw[brightsky!60!white, line width=5.5pt] (s6) -- (t6_9);
    \draw[brightsky!60!white, line width=5.5pt] (s6) -- (t6_10);
    \draw[vvgray, line width=11pt] (root) -- (s7);
    \draw[vvgray!60!white, line width=5.5pt] (s7) -- (t7_1);
    \draw[vvgray!60!white, line width=5.5pt] (s7) -- (t7_2);
    \draw[brightsky, line width=11pt] (root) -- (s8);
    \draw[brightsky!60!white, line width=5.5pt] (s8) -- (t8_1);
    \draw[brightsky!60!white, line width=5.5pt] (s8) -- (t8_2);
    \draw[vvteal, line width=11pt] (root) -- (s9);
    \draw[vvteal!60!white, line width=5.5pt] (s9) -- (t9_1);
    \draw[vvteal!60!white, line width=5.5pt] (s9) -- (t9_2);
    \draw[vvviolet, line width=11pt] (root) -- (s10);
    \draw[vvviolet!60!white, line width=5.5pt] (s10) -- (t10_1);
    \draw[vvviolet!60!white, line width=5.5pt] (s10) -- (t10_2);
    \draw[vvviolet!60!white, line width=5.5pt] (s10) -- (t10_3);
    \draw[vvviolet!60!white, line width=5.5pt] (s10) -- (t10_4);
    \draw[vvorange, line width=11pt] (root) -- (s11);
    \draw[tomato, line width=11pt] (root) -- (s12);
    \draw[tomato!60!white, line width=5.5pt] (s12) -- (t12_1);
    \draw[tomato!60!white, line width=5.5pt] (s12) -- (t12_2);
    \draw[tomato!60!white, line width=5.5pt] (s12) -- (t12_3);
  \end{scope}
\end{tikzpicture}
